\subsection{Optimización y función de costo}
\label{sec:funcion_de_costo}
En definitiva, obtener modelos que cumplan con los objetivos planteados, es equivalente a obtener modelos que optimizen estas métricas de evaluación.

En ciertos casos optimizar la métrica de evaluación es suficiente para obtener la funcionalidad deseada. Sin embargo, en la mayoría de los casos se quiere ponderar más de un aspecto del desempeño del modelo, incluso múltiples métricas de evaluación. Por ende, es común definir una función que combine estos diferentes aspectos en un solo valor cuantitativo que refleje el desempeño global del modelo. Esta función a optimizar es denominada función de costo $\mathcal{L}$.

Como se mencionó en la sección \ref{sec:modelado_de_sistemas} el problema fundamental que se busca resolver en el modelado de sistemas: encontrar los parámetros óptimos $\theta$ que minimicen la función de costo $\mathcal{L}$. Este proceso suele denominarse entrenamiento del modelo.

Existen múltiples técnicas de optimización empleadas para este tipo de problemas, \cite{Nocedal2006} describe una revisión exhaustiva de métodos de optimización numérica. En algunos casos, se dispone de soluciones cerradas que pueden calcularse analíticamente. Sin embargo, en la mayoría de los problemas se debe acudir a métodos numéricos iterativos para aproximar la solución óptima.

Ambos tipos de técnicas, tanto las soluciones cerradas como los métodos numéricos iterativos dependen de un conjunto de datos de entrada y salida etiquetados denominado conjunto de entrenamiento $\mathcal{D} = \left\{(x_i, y_i)\right\}_{i=1}^{N}$ (o \textit{dataset} en inglés).

En el caso del aprendizaje automático, el verdadero objetivo es minimizar la función de costo en datos no vistos previamente, es decir, minimizar el error de generalización del modelo. Sin embargo, dado que no se dispone de estos datos durante el proceso de entrenamiento, se utiliza el conjunto de entrenamiento como una aproximación para minimizar la función de costo.

Un desafío intrínseco de desconocer la naturaleza de la distribución de los datos, y de depender de un conjunto finito de entrada y salida, es obtener un modelo que sea capaz de generalizar la relación aprendida a datos no vistos previamente. Además, el hecho de que el conjunto de datos contiene ruido e incertidumbre, puede llevar a que el modelo se ajuste demasiado a las particularidades del conjunto de entrenamiento, perdiendo la capacidad de generalización. Este fenómeno es conocido como sobreajuste (\textit{overfitting} en inglés) \cite{Bishop2006}.

Para ello, es común dividir el conjunto de datos en subconjuntos de entrenamiento, validación y prueba, con el fin de evaluar el desempeño del modelo en datos no utilizados durante el proceso de ajuste de parámetros.

Como se mencionó anteriormente, la función de costo $\mathcal{L}$ raramente tiene un único término, el componente principal suele ser el término de pérdida $\mathcal{L}_{l}$, que mide la discrepancia entre las predicciones del modelo ($\hat{y}$) y los valores reales ($y$) del conjunto de entrenamiento. Sin embargo, como se mencionó, también se desea ponderar otros aspectos. Por ello, es muy común añadir términos de regularización a la función de costo $\mathcal{L}_{r}$, que penalizan ciertas características del modelo, como la complejidad o la magnitud de los parámetros \cite{Goodfellow2016}. La forma típica de la función de costo se expresa como \eqref{eq:cost_function}.

\begin{equation}
\label{eq:cost_function}
\mathcal{L}(\theta, \lambda, \mathcal{D}) =
    \underbrace{%
        \frac{1}{N} \sum_{i=1}^{N} \mathcal{L}_{l}(y_i, \hat{y}_i)
    }_{\text{Pérdida}}
    +
    \underbrace{%
        \lambda \cdot \mathcal{L}_{r}(\theta)\vphantom{\frac{1}{N} \sum_{i=1}^{N} \mathcal{L}_{l}(y_i, \hat{y}_i)}
    }_{\text{Regularización}}
\end{equation}

Donde $\lambda$ es un hiperparámetro que controla la importancia relativa del término de regularización en comparación con el término de pérdida. A diferencia de los parámetros $\theta$ del modelo que son entrenados durante el proceso de entrenamiento, un hiperparámetro es un valor de configuración externo que se fija antes de iniciar el proceso de entrenamiento y no se ajusta durante el mismo \cite{Goodfellow2016}.
